\documentclass[a4paper]{article}

\usepackage{graphicx}
\usepackage{amsmath,amssymb}
\usepackage{minted}
\usepackage{multirow}
\usepackage{float}
\usepackage{algorithm}
\usepackage{algpseudocode}
\usepackage{todonotes}
\usepackage{forest}
\usepackage{xstring}
\usepackage{xcolor}
\usepackage[most]{tcolorbox}
\usepackage{xparse}
\usepackage{natbib}

\usetikzlibrary{graphs,graphdrawing}
\usegdlibrary{layered}

% for external graphviz
\input{/home/donkarlo/Dropbox/repo/nd_language_ai_project/src/nd_language_ai/formal/computer/programming/tex/latex/code_clips/external_graphviz.tex}

% for tags
\input{/home/donkarlo/Dropbox/repo/nd_language_ai_project/src/nd_language_ai/formal/computer/programming/tex/latex/part/control_sequence/kind/semtag/semtag.tex}

% for links
\input{/home/donkarlo/Dropbox/repo/nd_language_ai_project/src/nd_language_ai/formal/computer/programming/tex/latex/code_clips/seehere.tex}

\title{Structure and large Persian editor cleanup}
\date{2026-09-05}
\author{Mohammad Rahmani}

\begin{document}
   \maketitle

   \section{Bug}
   The Structure panel could leave a stale second highlight because clearing stopped when the first structure item after the current source line was reached. Large Persian LaTeX documents also remained expensive to type in because a relatively large rich-text Source window was laid out by \texttt{QTextEdit}, and unchanged block formatting could be reapplied during ordinary edits.

   \section{Fix}
   Structure selection and highlighting are now owned directly by \texttt{LatexStructurePanel}. Every row is cleared before the single semantic target is painted, and ordinary prose edits in large documents no longer schedule a complete Structure parse when they cannot affect line mapping or structural commands. The separate selection-integrity monkey patch was removed, and the large-document runtime controller no longer installs duplicate Structure or line-number patches.

   \texttt{LargeLatexEditor} keeps the complete canonical document in \texttt{LatexSourceBuffer} while reducing the rich \texttt{QTextEdit} materialization window from 40,000 to 8,000 characters. Large-mode \texttt{TextEditor} also skips redundant block-format rewrites when direction and line height are already correct. Full-document idle publication was relaxed to reduce work while the user is actively writing.

   \section{Scope}
   The source document itself is unchanged. This change affects editor presentation, Structure synchronization, and large-document typing performance only.

   \bibliographystyle{plainnat}
   \bibliography{/home/donkarlo/Dropbox/repo/nd_shared_research/bibliography/bibliography.bib}
\end{document}
